% Generated draft: requires mathematical and source review.
\documentclass[11pt]{article}
\usepackage[margin=25mm]{geometry}
\usepackage{fontspec}
\setmainfont{Latin Modern Roman}
\usepackage{amsmath,amssymb,mathtools,unicode-math}
\setmathfont{Latin Modern Math}
\usepackage{enumitem,needspace,xurl,hyperref}
\setlength{\parindent}{0pt}
\setlength{\parskip}{5pt}
\setlength{\emergencystretch}{3em}
\newcommand{\N}{\mathbb{N}}
\newcommand{\Z}{\mathbb{Z}}
\newcommand{\Q}{\mathbb{Q}}
\newcommand{\R}{\mathbb{R}}
\newcommand{\C}{\mathbb{C}}
\newcommand{\F}{\mathbb{F}}
\newcommand{\A}{\mathbb{A}}
\newcommand{\PP}{\mathbb{P}}
\newcommand{\E}{\mathbb{E}}
\newcommand{\eps}{\varepsilon}
\newcommand{\dd}{\,\mathrm{d}}
\begin{document}
\section*{UnsolvedMath ID: 30001398}
\textbf{Source problem code:} catalogue entry
\subsection{Definitions and mathematical statement}
The supplied question concerns a hyperbolic exponential map \(f(z)=\lambda e^z\). Under the usual convention for nonconstant complex exponential maps, \(\lambda\in\C\setminus\{0\}\), and hyperbolicity means that the singular value \(0\) is attracted to an attracting periodic orbit. These conventions clarify terminology; the record does not explicitly state the parameter domain or define hyperbolicity.

For real numbers \(a<b\), write
\[
S_{a,b}=\{z\in\C:a\leq\operatorname{Im}z\leq b\}.
\]
This is one possible boundary convention for a horizontal strip of finite width. The supplied record does not specify whether strip boundaries are included.

A Hausdorff gauge is a nondecreasing function \(h:[0,\infty)\to[0,\infty)\) with \(h(0)=0\) and \(\lim_{r\downarrow0}h(r)=0\). With Euclidean diameter and an unnormalized convention, its Hausdorff measure is
\[
H^h(E)=\lim_{\delta\downarrow0}\inf\left\{\sum_{j=1}^{\infty}h(\operatorname{diam}U_j):E\subseteq\bigcup_{j=1}^{\infty}U_j,\ \operatorname{diam}U_j\leq\delta\right\}.
\]
The record supplies neither the gauge \(h\) nor the underlying set on which this measure is to be considered. In particular, it does not identify a Julia set, escaping set, or other carrier for the measure.

For clarity, one standard simultaneous Diophantine interpretation of extremality can be stated conditionally. Identify \(\C\) with \(\R^2\) by \(z\mapsto(\operatorname{Re}z,\operatorname{Im}z)\), and, for \(\varepsilon>0\), define
\[
W_\varepsilon=\left\{x\in\R^2:\|qx-p\|_\infty<q^{-1/2-\varepsilon}\text{ for infinitely many }(p,q)\in\Z^2\times\Z_{\geq1}\right\}.
\]
In this interpretation, a Borel measure \(\mu\) is extremal if \(\mu(W_\varepsilon)=0\) for every \(\varepsilon>0\).

If the intended gauge \(h=h_f\), a Borel carrier \(X_f\subseteq\C\), the preceding Diophantine interpretation, and the strip boundary convention are supplied, the question becomes: for every hyperbolic exponential map \(f\) and every \(a<b\), does the stipulated finite measure
\[
\mu_{f,a,b}(E)=H^{h_f}(E\cap X_f\cap S_{a,b}),\qquad E\subseteq\C\text{ Borel},
\]
satisfy
\[
\mu_{f,a,b}(W_\varepsilon)=0\qquad\text{for every }\varepsilon>0,
\]
where \(W_\varepsilon\) is transported to \(\C\) by the stated identification? This is a conditional formulation, not a resolution of the missing source specifications. The word ``finite'' is retained as a supplied qualification; positivity is not stated.

There is also an unresolved mismatch within the record: its statement asks about extremality, whereas its background and literature assessment describe Problem 10 as asking for a gauge representing a conformal measure. The record does not define that conformal measure or explain how the two questions are related.
\subsection{Short English statement}
The supplied question asks whether the finite Hausdorff measures associated with a hyperbolic exponential map, restricted to every horizontal strip of finite width, are extremal in the sense attributed to Kleinbock, Lindenstrauss, and Weiss. A conditional standard interpretation is given above, but the gauge, measure carrier, strip boundaries, and exact Diophantine interpretation require clarification. The accompanying assessment instead concerns representation of a conformal measure, so the intended target also requires review.
\subsection{Sources}
\begin{itemize}
\item Supplied local UnsolvedMath registry record, permanent ID \(30001398\), title: \emph{Extremal Hausdorff Measures for Hyperbolic Exponential Maps}; collection: Oberwolfach Reports Open Problems; exact source problem code: \texttt{catalogue entry}. No individual proposer is supplied; the proposed year is recorded as 2009.
\item Supplied authoritative reference: \emph{Mini-Workshop: The Escaping Set in Transcendental Dynamics} (2010), Oberwolfach Reports, pp. 2927--2964, DOI \texttt{10.4171/owr/2009/54}. Source URL: \url{https://doi.org/10.4171/owr/2009/54}. The literature-sources entry labels this same link ``original Oberwolfach problem report'': \url{https://doi.org/10.4171/owr/2009/54}.
\item Literature assessment dated exactly \texttt{2026-08-21}: the supplied curation labels the problem open and states that the report poses the existence of a gauge representing the conformal measure as Problem 10, with no later primary result located there establishing such a representation. This is supplied dated curation, not an independently verified assessment; it does not by itself establish the status of the extremality question.
\item The supplied statement-status label is \texttt{corrected\_verified}. Its accompanying verification note says that the mathematical content was checked against the official report and that notation, context, and question form were normalized without changing the source-level claim. This is a supplied verification claim, not verification performed for this draft.
\item Extremality is attributed in the supplied statements to Kleinbock, Lindenstrauss, and Weiss; the original-statement field names Barak Weiss. No separate bibliographic reference defining their notion of extremality is supplied.
\end{itemize}
\end{document}
